% Cote 119: offprints, reports, courses, CV and letters (typed and printed part): archive cote 119, site page 270 (PDF page 271).
% typescript. Draft compared: none.
% AI-transcribed by Claude Opus 5.5 (claude-opus-5-5, adaptive thinking, effort high), Message Batches API, batch r119_full (msgbatch_01DwZNm5Y7YsKBWyAN4sFs8w), 2026-11-13; not yet checked by eye.
\scanpage{119/270}
\mspage{17}

gnants. Quant aux étudiants,qui sont en principe la raison d'être de l'Univer-%
sité, leurs besoins dans tout cela ont été pratiquement ignorés. C'est là
un état de faits absurde et injuste, mais consacré par le passé, et d'une inertie
immense. Ceux qui détiennent toute l'autorité dans ce système, à savoir les
chercheurs, n'y renonceront pour la plupart sans une résistance extrêmement
sérieuse, passive ou active suivant le tempérament personnel ou les consignes
partisanes, suivant la conjoncture aussi. Dans un tel état d'esprit et faute
d'un vrai désir de faire du neuf, les discussions entre collègues et sans la
participation d'étudiants, sur tels ou tels points d'organisation et de
restructuration (comme ceux abordés dans les pages qui précèdent), risquent
évidemment de devenir byzantines, et même de donner finalement aux immobilistes
résolus une occasion de ``noyer le poisson''. Aussi il faut garder présent à
l'esprit que toutes les solutions proposées sont en dernière analyse des points
de détail, toujours discutables tant qu'elles ne sont pas sanctionnées et
corrigées à mesure par l'expérience, dans l'esprit de la ``contestation perma-%
nents'' sans lequel rien qui vaille ne saurait finalement être construit. Un
seul point de structure finalement est vraiment important, qui seul peut cons-%
tituer une garantie relative contre la tentation des maîtres de l'Université
d'enterrer au plus vite le mouvement : c'est la cogestion effective, et plus
précisément, la participation des étudiants, par une représentation massive
au sein des instances dirigeantes du département, à \uline{toutes} les décisions qui
les concernent. Ce point acquis, et les rapport\hand{s} de force au niveau des décisions \margin{s}
correspondant enfin aux rapports réels entre les fonctions des uns et des
autres, enseignants et étudiants, les structures s'ajusteront d'elles-mêmes ;
non en un jour ni en une année, mais chaque chose en son temps. C'est donc sur
le principe et les modalités d'une cogestion effective et dans l'immédiat que
tous ceux qui veulent que quelque chose change doivent, sous peine de s'épar-%
piller, concentrer leurs principaux efforts.

\margin{270}
